\section{Language of category}

\subsection{Functor and natural transformation}

\begin{definition}
    Let \(F: \mathcal{C} \to \mathcal{D}\) be a functor, we define
    \begin{itemize}
        \item \(F\) is \textbf{faithful} if for any two objects \(X,Y\) in \(\mathcal{C}\), the map
        \[F_{X,Y}: \Hom_{\mathcal{C}}(X,Y) \to \Hom_{\mathcal{D}}(F(X),F(Y))\]
        is injective.
        \item \(F\) is \textbf{full} if for any two objects \(X,Y\) in \(\mathcal{C}\), the map
        \[      F_{X,Y}: \Hom_{\mathcal{C}}(X,Y) \to \Hom_{\mathcal{D}}(F(X),F(Y))\]
        is surjective.      
        \item \(F\) is faithful and full if such map is bijective.
        \item \(F\) is \textbf{essentially surjective} if for any object \(Y\) in \(\mathcal{D}\), there exists an object \(X\) in \(\mathcal{C}\) such that \(F(X)\) is isomorphic to \(Y\); Furthermore, if \(F(X) = Y\) we say \(F\) is surjective.
    \end{itemize}
\end{definition}

\begin{remark}
    Such a definition can not be seen as a "injection/surjection between category" in the level of objects.
\end{remark}

\begin{example}[full but not faithful]
    Consider the functor from \(\mathbf{Top}\) to \(\mathbf{Homotopy}\), in the level of morphism, it maps a continuous map \(f: X \to Y\) to its homotopy class \([f]: [X] \to [Y]\). 
    
\end{example}

\begin{example}[essentially surjective]
    Let \(\mathcal{C}_{\nn}\) be the discrete category of natural numbers, then
    \[\mathbf{Vect}^{f,B}_{k} \to \mathcal{C}_{\nn}, (V,B) \mapsto \dim(V)\]
    is surjective, by comparison
    \[\mathcal{C}_{\nn} \to \mathbf{Vect}, n \mapsto (k^n,standard basis)\]
    is not surjective by considering \(V = k[X]/(X^n)\), but it is essentially surjective since \(V \cong k^n\).
\end{example}

\begin{example}
    Consider the discrete category \(\mathcal{C} \) of \(\pi_0(X)\), i.e. the set of path-connected components of a topological space \(X\), 
    \[a \mapsto [a]_0 = \{\text{ comonent of a in X}\}\]
    gives a surjective functor from groupoid \(\pi_{\leq 1}(X)\) to above category, conversely we consider 
    \[C \in \pi_0(X) \mapsto \text{some point of component}\]
    it gives a functor which is not surjective but essentially surjective.
\end{example}

\begin{colordefinition}[natural transformation]
    Let \(F,G: \mathcal{C} \to \mathcal{D}\) be two functors, a natural transformation \(\eta: F \to G\) is a collection of \(\mathcal{D}\)-morphisms \(\eta_X: F(X) \to G(X)\) for each object \(X\) in \(\mathcal{C}\) such that for any morphism \(f: X \to Y\) in \(\mathcal{C}\), the following diagram commutes:
    \[
    \begin{tikzcd}
    F(X) \arrow[r, "\eta_X"] \arrow[d, "F(f)"'] & G(X) \arrow[d, "G(f)"] \\
    F(Y) \arrow[r, "\eta_Y"'] & G(Y)
    \end{tikzcd}
    \]
    sometimes we write \(F \Rightarrow G\) to denote the natural transformation.
    
\end{colordefinition}

\begin{example}
    For any category \(\mathcal{C}\), we denote \(I_{\mathcal{C}}\)
to be the identity functor on \(\mathcal{C}\), in the case of vector space of dimension, bidual gives a natural transformation
    \[\eta: I_{\mathbf{Vect}_k} \to (-)^{**}\]
    where \(\eta_V: V \to V^{**}\) is the canonical map \(v \mapsto (f \mapsto f(v))\).
\end{example}

\begin{example}
    \(GL_n(-)\) and \(\det (-)\) in commutative ring.
\end{example}

\textcolor{blue}{Natural transformation allows us to construct a 2-category, i.e. the generalization of category in some sense.}

\begin{definition}[category of functors]
    Let \(\mathcal{C}\) and \(\mathcal{D}\) be two catgories, we define
    \[\cat{Fun(\mathcal{C},\mathcal{D})}{\text{all functors from }\mathcal{C}\text{ to }\mathcal{D}}{\text{natrual tranfromations}}\]
\end{definition}

\begin{example}
    Functor \(\Hom(-,X)\) and \(\Hom(X,-)\) can be desribed by natural transformation, as the target of the morphism of 2-cat.
\end{example}

\begin{remark}
    For sake of convience, we make a convention that for any category \(\mathcal{C}\), we write
    \[\mathcal{C}^\vee = Fun(\mathcal{C}^{\text{op}}, \mathbf{Set})\]
    while
    \[\mathcal{C}^{\wedge} = Fun(\mathcal{C}^{\text{op}}, \mathbf{Set})^{\text{op}}\]
    and \(\mathcal{C}^\vee\) is called presheaf on \(\mathcal{C}\). 
\end{remark}

\begin{example}
    Let \(X\) be a topological space with \(\mathcal{C}\) the poset category of open sets of \(X\), then \(\mathcal{C}^\vee\) is the category of sheaves on \(X\).
\end{example}

\begin{colordefinition}
    The functor \(F: \mathcal{C} \to \mathcal{D}\) is said to be an \textbf{isomorphism of catgeories} if \(F\) is faithful, full and bijective on objects.\\

    \(F\) is said to be an \textbf{equivalence of catgeories} if \(F\) is faithful, full and essentially surjective.\\
 \end{colordefinition}

 \begin{example}
    In the context of representation theory, we have an isomorphism of catgeories
    \[\mathbf{Rep}_k(G) \cong k[G]-\mathbf{Mod}\]
    where \(G\) is a group and \(k\) is a field.
 \end{example}

 \begin{example}
    Let \(X\) be a locally path connected topological space, then 
    \[F: \pi_{\leq 1}(X) \to \mathcal{C}_{\pi_0(X)}\]
    is not an isomoprhism of categories since it is not injective on objects; and it is an equivalence if each compoenent of \(X\) is simply connected.
 \end{example}

 \begin{example}
    A category \(\mathcal{C}\) is said to be discrete if there exists a set \(E\) such that \(\mathcal{C} \cong \mathcal{C}_{E}\) by equivalence, by previous example, \(\mathbf{Vect}_k^f\) is discrete.
 \end{example}

\begin{colordefinition}
    Let \(F,G: \mathcal{C} \to \mathcal{D}\) be two functors, an \textbf{isomorphism of functors} \(\eta: F \to G\) is a natural transformation such that for each object \(X\) in \(\mathcal{C}\), the morphism \(\eta_X: F(X) \to G(X)\) is an isomorphism in \(\mathcal{D}\).\\

    Or equivalently, there exists a natural transformation \(\beta: G \to F\) such that \(\beta \circ \eta = I_F\) and \(\eta \circ \beta = I_G\).\\

    Or equivalently, \(\eta\) is an isomorphism as morphism in the 2-cat \(Fun(\mathcal{C},\mathcal{D})\).
\end{colordefinition}


\subsection{Adjoint}

\begin{definition}
    IF $F: \mathcal{C} \to \mathcal{D}$ and $G: \mathcal{D} \to \mathcal{C}$ are two covariant functors, then \((F,G)\) is pair of adjoint functors if there is a collection of bijections \(u_{X,Y}\) for all ojects \(X\) in \(\mathcal{C}\) and \(Y\) in \(\mathcal{D}\) such that
    \[u_{X,Y}: \Hom_{\mathcal{D}}(F(X), Y) \cong \Hom_{\mathcal{C}}(X, G(Y))\] 
\end{definition}

\begin{remark}
    such a definition can be rephrased in terms of product of catgeory.
\end{remark}

For notation we write \(F \dashv G\) to denote that \(F\) is left adjoint to \(G\) or \(G\) is right adjoint to \(F\).

\begin{example}
    In the context of Kaler module, we have a universal property 
    \[\Hom_B(\Omega_{B/A},M) \cong \mathrm{Der}_A(B,M)\]
    where \(B\) is a \(A\)-algebra, \(M\) is a \(B\)-module, \(\Omega_{B/A}\) is the module of Kähler differentials and \(\mathrm{Der}_A(B,M)\) is the set of \(A\)-derivations from \(B\) to \(M\). Such property can be described in 
    \[\cat{A-alg/C}{}{}\]
\end{example}